\section{Background}\label{sec:background}

\subsection{Anatomy of an LLM Decode Step}

Autoregressive decoding generates one or a few tokens per request per step.
Unlike prefill, where matrix multiplications operate on long prompts and are compute-bound,
low-batch decode has low arithmetic intensity: every step streams the model's full weight set
from off-chip memory for a handful of token rows, making decode memory-bandwidth bound
under the roofline model~\cite{williams2009roofline}.
Each transformer or hybrid layer issues a sequence of weight-streaming GEMMs
(input projections, output projection, and MLP projections), a sequence-mixing kernel
(attention or a recurrent state update), and small glue kernels (RMSNorm, RoPE,
activation quantization, residual adds).

To avoid host-side launch latency across hundreds of short kernels per step,
serving engines such as vLLM~\cite{kwon2023vllm} capture steady-state decode steps into
CUDA graphs and replay them as a single submission.
Inside a replayed graph, kernels execute back-to-back on the GPU with minimal inter-kernel
idle time; performance is governed by how many bytes each launch moves and how close its
effective bandwidth comes to the device's achievable memory bandwidth ceiling.
Because a single compiled GEMM kernel template (for example, a CUTLASS FP8 kernel)
is reused across dozens of weights in different layers, aggregating execution time by kernel
name collapses distinct weights and distinct cache contexts into one number.

\subsection{Speculative Decoding with Multi-Token Prediction}

Speculative decoding~\cite{leviathan2023fast} raises arithmetic intensity during decode by
verifying $k$ draft tokens in parallel alongside the current token, passing $k{+}1$ tokens
per request through the main model in a single step.
Because weight reads are amortized across all $k{+}1$ tokens, verification costs nearly the
same memory traffic as a single-token step while yielding up to $k{+}1$ accepted tokens.

Multi-Token Prediction (MTP)~\cite{deepseekv3} integrates drafting into the model itself via
lightweight prediction heads that run after the main forward pass.
With $k{=}3$ draft tokens, each decode step runs the main model on four tokens and then
invokes the draft head three times sequentially, reading the vocabulary projection
(\texttt{lm\_head}) once for the main model and once per draft step.

\subsection{Hybrid Gated DeltaNet Models and State Rollback}

Hybrid architectures interleave standard softmax attention layers with linear-time recurrent
layers to reduce KV-cache growth at long contexts.
Gated DeltaNet (GDN)~\cite{yang2025gated} maintains a fixed-size matrix-valued recurrent
state per head---in Qwen3.8-27B, 48 heads of $128 \times 128$ in FP32, or 3.0~MiB per
sequence per layer---across 48 of the model's 64 layers (with 16 standard attention layers).

Speculative verification complicates in-place recurrent updates: if the verifier rejects
draft token $j \in \{1, \dots, k\}$, the recurrent state must roll back to position $j{-}1$
before the next step.
The straightforward implementation writes the full updated state after every verified
position into a dedicated slot.
For $k{=}3$ draft tokens (four positions per step), a single GDN decode launch writes
$4 \times 3.0 = 12$~MiB of recurrent state per request per layer, even though only the
single slot at the last accepted position is ever read on the next step.
When these stores land in a write-back L2 cache, the GDN kernel can retire quickly while
leaving megabytes of dirty cache lines whose DRAM write-back traffic contends with the
weight-streaming GEMM that immediately follows.

\subsection{Target Platform: NVIDIA DGX Spark (GB10)}

Our case study runs on an NVIDIA DGX Spark built around the GB10 SoC
(\texttt{sm\_121}, 48 SMs, 24~MiB L2 cache) backed by 121~GB of unified LPDDR5x memory.
While the datasheet quotes 273~GB/s peak bandwidth, microbenchmarks---our read-only
streaming kernel (\texttt{bwtest}) and BabelStream~\cite{deakin2016babelstream}---consistently
measure 248--253.5~GB/s (we use 252~GB/s as the practical ceiling).
GB10 operates under a tight software power cap during sustained decode, causing whole-step
wall-clock times to drift by 1--2~ms across runs even with the GPU clock locked at 2000~MHz.
Reliable attribution on such hardware therefore requires per-kernel measurements paired with
control kernels (unchanged layers and projections) rather than whole-step averages alone.
